\documentclass{article}
\usepackage{amsmath} % Required for advanced math features like 'align'
\title{A Sample Document with Mathematics}
\author{My Name}
\date{\today}

\begin{document}

\maketitle

\section{Introduction}
\LaTeX{} is excellent for writing documents with mathematical notation. Mathematical expressions can be included inline within a paragraph using single dollar signs, like this: the famous Pythagorean theorem is often written as $a^2 + b^2 = c^2$.

\section{Displayed Equations}
For equations that need to be on their own line and possibly numbered, the `equation` environment is ideal.

\begin{equation}
E = mc^2 \label{eq:einstein}
\end{equation}

Equation \ref{eq:einstein} is the mass-energy equivalence principle.

For unnumbered displayed equations, you can use double dollar signs (`$$...$$`) or `\[...\]` (the latter is preferred).

\[
\sum_{i=1}^{N} i = \frac{N(N+1)}{2}
\]

\section{Advanced Mathematics}
The `amsmath` package provides environments for more complex or aligned equations, such as the `align*` environment (the asterisk prevents numbering).

\begin{align*}
f(x) &= (x+1)^2 \\
     &= x^2 + 2x + 1
\end{align*}

We can also include matrices using environments like `pmatrix` or `array`.

\[
A = \begin{pmatrix}
a_{11} & a_{12} \\
a_{21} & a_{22}
\end{pmatrix}
\]

\end{document}